\documentclass[a4paper, 11pt]{article}

\usepackage[margin=2cm]{geometry}
\usepackage{booktabs}
\usepackage{array}
\usepackage{tabularx}
\usepackage{amsmath}
\usepackage{amssymb}
\usepackage{caption}
\usepackage{hyperref}
\usepackage[T1]{fontenc}
\usepackage{lmodern}

% Column types
\newcolumntype{L}[1]{>{\raggedright\arraybackslash}p{#1}}
\newcolumntype{C}[1]{>{\centering\arraybackslash}p{#1}}

\title{Parameter Identification Experiments\\
\large Garc\'{i}a et al.\ (2013) --- Dual-Drive Gantry Stage}
\author{}
\date{}

\begin{document}
\maketitle

\section*{Source}
I.\ Garc\'{i}a-Herreros, R.\ Coleman, P.-J.\ Barre, J.\ Gomand, X.\ Kestelyn,
\textit{Model-based decoupling control method for dual-drive gantry stages:
A case study with experimental validations},
Control Engineering Practice, Vol.\ 21, No.\ 3, pp.\ 298--307, 2013.

The identification procedure is described in \textbf{Section 2.4} (``Model parameters'')
and is organised as a sequential chain: each experiment assumes all prior parameters
are already known.

\section*{Model parameters (14 physical parameters)}

The lumped-parameter model (Section~2.2) is described by fourteen physical parameters:

\begin{itemize}
    \item \textbf{4 masses:} $m_b$ (cross-arm), $m_h$ (payload/Y-axis), $m_1$ (actuator $X_1$), $m_2$ (actuator $X_2$)
    \item \textbf{5 viscous friction coefficients:} $c_{g1}$, $c_{g2}$ ($X_1$, $X_2$), $c_y$ (Y-axis), $c_{b1}$, $c_{b2}$ (flexible joints)
    \item \textbf{3 Coulomb friction forces:} $c_{c1}$, $c_{c2}$ ($X_1$, $X_2$), $c_{cy}$ (Y-axis)
    \item \textbf{2 stiffness coefficients:} $k_{b1}$, $k_{b2}$ (flexible joints)
\end{itemize}

\noindent Note: the rotary inertias $J_b$ and $J_h$ appear in the model but are not free
parameters --- they are resolved via the bar assumption Eq.~(21):
$J_b = m_b L_b^2 / 12$, $J_h = m_h L_h^2 / 12$.

\section*{Identification experiments}

\begin{table}[htbp]
\centering
\caption{Sequential parameter identification experiments (Section~2.4, Garc\'{i}a et al.\ 2013).
Experiments must be executed in order: each step uses results from all previous steps.}
\label{tab:garcia_experiments}
\small
\renewcommand{\arraystretch}{1.6}
\begin{tabular}{C{0.5cm} C{0.9cm} L{3.2cm} C{1.0cm} L{3.2cm} L{5.0cm}}
\toprule
\textbf{Exp} & \textbf{Axis} & \textbf{Motion condition} & \textbf{Eq.} & \textbf{Parameters identified} & \textbf{Physical mechanism} \\
\midrule

1 & Y &
$\dot{Y} = \text{const}$, $\ddot{Y} = 0$;\newline $X_1 = X_2 = \text{const}$ (blocked) &
(17) &
$c_y$,\ $c_{cy}$ &
Eq.~(17) reduces to $F_y = c_y \dot{Y} + c_{cy}\,\mathrm{sign}(\dot{Y})$. Sweep multiple
constant velocities: slope $\rightarrow c_y$, intercept $\rightarrow c_{cy}$. \\

\midrule

2 & Y &
$\ddot{Y} = \text{const}$;\newline $X_1 = X_2 = \text{const}$ (blocked) &
(17) &
$m_h$,\ $d$ &
Known $c_y$, $c_{cy}$ from Exp.~1. $m_h\ddot{Y} = F_y - c_y\dot{Y} - c_{cy}\,\mathrm{sign}(\dot{Y})$
gives $m_h$. Asymmetric reaction $F_1 - F_2 = m_h d\,\ddot{Y}$ gives the CoM offset $d$. \\

\midrule

3 & $X_1, X_2$ &
$\dot{X}_1 = \dot{X}_2 = \text{const}$,\newline $\ddot{X}_1 = \ddot{X}_2 = 0$ &
(18),(19) &
$c_{g1}$,\ $c_{g2}$,\newline $c_{c1}$,\ $c_{c2}$ &
$F_1 + F_2$ gives $(c_{g1}+c_{g2})\dot{X} + (c_{c1}+c_{c2})\mathrm{sign}(\dot{X})$;
$F_1 - F_2$ gives the differences. Fitting at multiple velocities separates all four values. \\

\midrule

4 & $X_1, X_2$ &
$\ddot{X}_1 = \ddot{X}_2 = \text{const}$ &
(18),(19) &
$m_1{+}m_2{+}m_b{+}m_h$,\newline $m_1{-}m_2$ &
Frictions known from Exp.~3. Eq.~(18) $\rightarrow$ total translational mass.
Eq.~(19) $\rightarrow$ mass asymmetry $m_1 - m_2$.
Individual $m_b$, $m_1$, $m_2$ remain underdetermined until Exp.~7 + Eq.~(21). \\

\midrule

5 & $\Theta$ &
$\Theta$ held at fixed angles\newline $-0.1, -0.09, \ldots, +0.1$\,rad;\newline $\dot{\Theta} = \ddot{\Theta} = 0$ &
(14) &
$k_{b1}{+}k_{b2}$ &
Static equilibrium: $(k_{b1}+k_{b2})\Theta = (F_1-F_2)\,L_b/2$.
Linear fit across multiple fixed angles gives the stiffness sum. \\

\midrule

6 & $\Theta$ &
$\dot{X}_1 = -\dot{X}_2 = \text{const}$\newline $\Rightarrow\ \dot{\Theta} = \text{const}$, $\ddot{\Theta} = 0$ &
(13) &
$c_{b1}{+}c_{b2}$ &
No inertia term. Rotational damping: $[c_{b1}+c_{b2}+(c_{g1}+c_{g2})L_b^2/4]\,\dot{\Theta}$
equals measured torque. $c_{g1}+c_{g2}$ known from Exp.~3, so $c_{b1}+c_{b2}$ is extracted. \\

\midrule

7 & $\Theta$ &
$\ddot{X}_1 = -\ddot{X}_2 = \text{const}$\newline $\Rightarrow\ \ddot{\Theta} = \text{const}$ &
(20),(21) &
$J_b{+}J_h{+}(m_1{+}m_2)\tfrac{L_b^2}{4}{+}m_h d^2$;\newline resolves $m_b$, $m_1$, $m_2$ &
All friction and stiffness known. Eq.~(20) gives the total rotary inertia sum.
Eq.~(21) ($J_b = m_b L_b^2/12$, $J_h = m_h L_h^2/12$) closes the algebraic system
to separately recover $m_b$, $m_1$, $m_2$. \\

\bottomrule
\end{tabular}
\end{table}

\section*{Dependency chain}

\[
\underbrace{\text{Exp.~1} \to \text{Exp.~2}}_{\text{Y-axis (decoupled)}}
\;\longrightarrow\;
\underbrace{\text{Exp.~3} \to \text{Exp.~4}}_{\text{Synchronized X}}
\;\longrightarrow\;
\underbrace{\text{Exp.~5} \to \text{Exp.~6} \to \text{Exp.~7} \xrightarrow{\text{Eq.~(21)}} m_b,\,m_1,\,m_2}_{\text{Rotational (uses all prior)}}
\]

% ---------------------------------------------------------------
\newpage
\section*{LPV-LFR baseline trajectories (T1--T6)}

The six trajectories are generated by \texttt{export\_lpv\_multi\_traj.m} using
third-order (jerk-limited) set-point profiles. Unlike Garc\'{i}a, all trajectories
are fitted \emph{simultaneously} via gradient-based BPTT; there is no sequential
algebraic chain. Identification uses position output only --- force measurements
are not available.

\medskip
\noindent\textbf{Key structural differences from Garc\'{i}a:}
\begin{itemize}
    \item Profile shape: smooth third-order (jerk-limited) vs.\ piecewise constant velocity/acceleration.
    \item Coulomb friction ($c_{c1}$, $c_{c2}$, $c_{cy}$): \textbf{not modelled} in the LFR --- absent from all trajectories.
    \item CoM offset $d$ and cross-arm length $L_b$: \textbf{fixed buffers} (0.1\,m, 0.725\,m) --- not trainable.
    \item T5 (simultaneous X+Y) is \textbf{LPV-unique}: exploits the scheduling variable $Y$ varying
          continuously; Garc\'{i}a has no equivalent.
    \item No static-hold equivalent of Garc\'{i}a Exp.~5 (fixed $\Theta$ angles for stiffness).
\end{itemize}

\begin{table}[htbp]
\centering
\caption{LPV-LFR baseline trajectories T1--T6 and their relationship to
Garc\'{i}a's identification experiments.
$v$/$a$ limits are given as (X, Y) pairs where both axes move; Y-only or X-only
trajectories list the active axis only.}
\label{tab:trajectories}
\small
\renewcommand{\arraystretch}{1.6}
\begin{tabular}{C{0.5cm} L{2.8cm} L{3.0cm} L{3.0cm} C{1.1cm} L{3.5cm}}
\toprule
\textbf{Traj} & \textbf{Motion} & \textbf{Parameters targeted} & \textbf{Physical mechanism} & \textbf{García Exp.} & \textbf{Comparison} \\
\midrule

T1 &
Y sweep, conservative\newline $Y: 0.3\to-0.3$\,m\newline $v_Y=1.0$\,m/s,\ $a_Y=10$\,m/s$^2$ &
$c_y$,\ $m_h$ &
At lower dynamics the friction term $c_y\dot{Y}$ and inertia term $m_h\ddot{Y}$
are comparably sized, exciting both simultaneously. &
Exp.~1, 2 &
\textbf{Partial.} Covers the same Y-only excitation as Exp.~1+2 combined.
Smooth profile instead of piecewise constant; Coulomb $c_{cy}$ not modelled. \\

\midrule

T6 &
Y sweep, aggressive\newline $Y: 0.3\to-0.3$\,m\newline $v_Y=2.0$\,m/s,\ $a_Y=50$\,m/s$^2$ &
$m_h$ (primary),\ $c_y$ (secondary) &
Higher dynamics shift dominance toward $m_h\ddot{Y}$.
T1+T6 together: $c_y v/(m_h a)$ ratio = 0.099 (T1) vs.\ 0.040 (T6),
allowing the optimizer to disambiguate $c_y$ from $m_h$. &
Exp.~1, 2 &
\textbf{Partial.} T6 at hardware-maximum dynamics is the closest analog to
Garc\'{i}a's constant-acceleration mass identification. Still a smooth profile. \\

\midrule

T2 &
X symmetric at $Y=0.3$\,m\newline $X_1{=}X_2: 0\to150$\,mm\newline $v_X=1.5$\,m/s,\ $a_X=20$\,m/s$^2$ &
$c_{g1}$,\ $c_{g2}$,\ $m_1{+}m_2{+}m_b$,\ $m_h$\ (via $M_1$ coupling) &
Synchronized X at $Y=0.3$\,m maximises LPV coupling:
$M_1[0,1] = -m_h Y = -3.03$\,kg$\cdot$m.
Friction excited by $\dot{X}$; inertia by $\ddot{X}$. &
Exp.~3, 4 &
\textbf{Partial.} Same synchronized-X mode as Garc\'{i}a. LPV coupling at $Y=0.3$\,m
has no Garc\'{i}a analog. Coulomb $c_{c1}$, $c_{c2}$ not modelled. \\

\midrule

T3 &
X symmetric at $Y=0$\,m\newline $X_1{=}X_2: 0\to150$\,mm\newline $v_X=1.5$\,m/s,\ $a_X=20$\,m/s$^2$ &
$c_{g1}$,\ $c_{g2}$,\ $m_1{+}m_2{+}m_b$ &
Identical motion to T2 but $M_1$ coupling $= 0$ at $Y=0$.
T2 vs.\ T3 contrast isolates $m_h$ from $m_1{+}m_2{+}m_b$. &
Exp.~3, 4 &
\textbf{Partial.} T2+T3 together achieve the same mass disambiguation as
Exp.~3+4, but via a Y-position contrast rather than a sequential algebraic formula. \\

\midrule

T4 &
X anti-symmetric at $Y=0.2$\,m\newline $X_1{=}+35$\,mm,\ $X_2{=}-35$\,mm\newline $v_X=0.5$\,m/s,\ $a_X=8$\,m/s$^2$ &
$k_{b1}{+}k_{b2}$,\ $c_{b1}{+}c_{b2}$,\ $J_b{+}J_h$ &
Only trajectory exciting the rotational mode ($\dot{X}_1=-\dot{X}_2$).
Stiffness, joint damping, and rotary inertia all appear in the $\Theta$-equation.
Low dynamics chosen to keep $\Theta \leq 0.097$\,rad within safe limits. &
Exp.~5, 6, 7 &
\textbf{Weakest match.} Covers all three of Garc\'{i}a's rotational experiments
in one dynamic trajectory. Garc\'{i}a separates stiffness (static hold),
damping (const speed) and inertia (const accel) sequentially;
here all three are excited simultaneously in a single smooth move. \\

\midrule

T5 &
X sym.\ + Y sweep simultaneously\newline $X_1{=}X_2: 0\to100$\,mm,\newline $Y: 0.2\to-0.2$\,m\newline $v_X=1.0$, $a_X=15$\,m/s$^2$ &
$m_h$,\ $M_1(Y) = -m_h Y$ over full range &
Y varies continuously while X accelerates, tracing the full LPV mass-coupling
schedule. Provides gradient signal for $m_h$ across the entire operating range
of the scheduling variable. &
\emph{None} &
\textbf{LPV-unique.} Garc\'{i}a's model has no scheduling variable; this trajectory
has no equivalent in his identification procedure. \\

\bottomrule
\end{tabular}
\end{table}

\end{document}
